\section{Objectifs}
\begin{itemize}
\item Croiser deux caractères et choisir une représentation adaptée.
\item Lire un nuage de points et calculer un point moyen.
\item Utiliser un ajustement affine pour interpoler ou extrapoler.
\item Interpréter les résultats avec esprit critique.
\end{itemize}

\section{Prérequis}
Statistiques à une variable, fonctions affines, lecture de graphiques et usage élémentaire du tableur.

\section{1. Deux caractères : tableau croisé}
Une étude bivariée porte sur deux caractères observés sur les mêmes individus. Lorsque les deux caractères sont qualitatifs, un \\textbf{tableau croisé d'effectifs} répartit les individus selon les deux critères.

Exemple : dans un groupe A de 72 personnes, 30 possèdent aussi la caractéristique B. La fréquence de B parmi A vaut
\[
\frac{30}{72}\approx 0{,}417,
\]
soit environ 41,7 \\%.

Cette fréquence décrit le sous-groupe A ; elle ne doit pas être confondue avec la fréquence de B dans toute la population.

\section{2. Deux caractères quantitatifs : nuage de points}
Pour des couples $(x_i,y_i)$, on représente les données par un \\textbf{nuage de points}. Sa forme aide à repérer une tendance, sans démontrer à elle seule une relation de cause à effet.

Le \\textbf{point moyen} $G$ a pour coordonnées
\[
G(\bar x,\bar y),\qquad
\bar x=\frac{x_1+\cdots+x_n}{n},\qquad
\bar y=\frac{y_1+\cdots+y_n}{n}.
\]

\section{3. Ajustement affine}
Si le nuage présente une tendance approximativement linéaire, on peut choisir une droite d'ajustement
\[
y=ax+b.
\]

Le programme demande d'utiliser et d'interpréter un ajustement affine, sans imposer l'étude théorique de la méthode des moindres carrés.

\begin{itemize}
\item \\textbf{Interpoler} : estimer une valeur à l'intérieur de la zone couverte par les données.
\item \\textbf{Extrapoler} : estimer une valeur au-delà de cette zone.
\end{itemize}

Une extrapolation demande davantage de prudence : le modèle peut cesser d'être pertinent.

\section{4. Tableur et esprit critique}
Un tableur aide à trier, filtrer, représenter les données et tester un modèle. Le calcul automatique ne remplace pas l'interprétation : axes, unités, période étudiée et ordre de grandeur doivent toujours être vérifiés.

\section{Automatismes à entretenir}
\begin{itemize}
\item Lire correctement axes, unités et échelles.
\item Calculer moyenne et pourcentages simples.
\item Passer d'un tableau à une représentation graphique.
\item Déterminer le coefficient directeur d'une droite à partir de deux points.
\end{itemize}

\section{Exemple guidé}
Pour les points $(1,2)$, $(3,4)$ et $(5,6)$ :
\[
\bar x=3,\qquad \bar y=4.
\]
Le point moyen est donc $G(3,4)$.

Si une droite d'ajustement est $y=1{,}8x+12$, alors pour $x=25$ :
\[
y=1{,}8\times25+12=57.
\]

\section{Erreurs fréquentes}
\begin{itemize}
\item Confondre corrélation et causalité.
\item Oublier les unités.
\item Extrapoler très loin des données sans discuter la validité du modèle.
\item Lire un graphique sans vérifier l'origine ou l'échelle des axes.
\end{itemize}

\section{À retenir}
Le tableau croisé convient aux données qualitatives. Pour deux caractères quantitatifs, le nuage de points permet d'étudier une relation éventuelle et un ajustement affine peut fournir des estimations, à condition d'en discuter la pertinence.
